\documentclass[conference]{IEEEtran}
\usepackage{amsmath,amssymb,amsthm}
\usepackage[T1]{fontenc}
\usepackage{booktabs}
\usepackage{graphicx}
\usepackage{hyperref}
\hypersetup{hidelinks}
\newtheorem{theorem}{Theorem}
\newtheorem{lemma}{Lemma}
\newtheorem{corollary}{Corollary}
\newcommand{\Oh}{\mathcal{O}}
\title{Elementary Mixing Bounds and Experiments for the\\
Representation Technique in Subset Sum}
\author{\IEEEauthorblockN{Le Duc Minh}
\IEEEauthorblockA{\href{mailto:minh.leduc.0210@gmail.com}{minh.leduc.0210@gmail.com}}}
\begin{document}
\maketitle
\begin{abstract}
Subset sum is solved on every input in about $2^{n/2}$ steps, and nobody knows
whether $2^{n/3}$ (or anything below $2^{n/2}$) is possible. The representation
technique does much better on random inputs, about $2^{0.291n}$, and the missing
ingredient for general inputs is usually called \emph{mixing}: the many
representations of a solution must fall into many residue classes modulo a
random prime. This note does not give a faster worst-case algorithm. It gives
two elementary facts about mixing and one experiment. First, if the weight-$w$
sums take few residues modulo a prime $p$, then the input values themselves are
concentrated modulo $p$. Second, averaged over a window of primes, the collision
energy of the representations is at most $E_\infty+Y^2\kappa/\pi$, which is
$\Oh(Y\log Y)$ when few exact collisions occur, so a random prime mixes to within
a polylogarithmic factor on \emph{every} input. We then run the complete
pipeline---modular filter, balanced sub-solver, disjointness match---on
adversarial families. Low mixing appears only on arithmetic structure, which is
easy; the hard large-value and generalized-progression families mix at
$0.85$--$0.95$ and are solved at the sub-solver cost
$2\binom{n/2}{n/8}=2^{0.4057n}$. We conclude that for the coefficient set
$\{0,1\}$ the obstruction is more likely the balanced sub-solver than mixing, and
we state the resulting open problem. The two lemmas are machine-checked in Lean.
\end{abstract}
\begin{IEEEkeywords}
subset sum, representation technique, mixing, meet-in-the-middle, worst-case
analysis
\end{IEEEkeywords}

\section{Introduction}
Suppose you have $n$ integers and a target. Can you pick some of them so that the
picked numbers add up exactly to the target? This is the \emph{subset sum}
problem. For the numbers $3,5,8,13$ and target $16$ the answer is yes, in two
ways: $3+13$ and $3+5+8$.

Written as an equation, given integers $a_1,\dots,a_n$ and $b$, we ask whether
$\exists x\in\{0,1\}^n:\sum_i a_ix_i=b$.

\subsection{Why it matters}
Subset sum is one of the standard NP-complete problems. Trying every subset costs
$2^n$ steps; meet-in-the-middle of Horowitz and Sahni~\cite{hs} brings this to
about $2^{n/2}$, and that is still the best bound that holds on every input,
fifty years later. Schroeppel and Shamir~\cite{ss} kept the $2^{n/2}$ time with
$2^{n/4}$ memory, and Chen, Jin, Randolph and Servedio~\cite{chen} removed a
polynomial factor, but the exponent stayed $1/2$. Whether $2^{(1/2-\varepsilon)n}$
is possible is open.

\subsection{What is known}
On \emph{random} inputs the representation technique of Howgrave-Graham and
Joux~\cite{hgj}, refined by Becker, Coron and Joux~\cite{bcj}, runs in about
$2^{0.291n}$ steps. Its speed comes from splitting the \emph{answer} rather than
the positions: a solution has enormously many two-piece \emph{representations},
and after filtering them with a random prime only a small subfamily remains to be
searched. The analysis assumes the representations are spread across the residue
classes---they \emph{mix}. Randolph and W\k{e}grzycki~\cite{rw} made the
technique work on every input for many coefficient sets, using a ``mixing
dichotomy''; the plain subset sum case $\{0,1\}$ resisted. Austrin, Kaski,
Koivisto and Nederlof~\cite{akkn} showed that instances with a very small or very
large maximum bin size are also easy, leaving a hard middle band.

\subsection{What this note does}
\begin{itemize}
\item It sets up the mixing condition and defines the coverage and the collision
energy (Section~\ref{sec:mix}), with a worked example.
\item It proves two elementary lemmas: few hit residues force the input values to
be concentrated modulo the prime, and one residue class can be stripped and the
instance recursed (Section~\ref{sec:lemmas}). Both are machine-checked in Lean.
\item It proves that the average collision energy over a window of primes is
$\Oh(Y\log Y)$ when few exact collisions occur, so a random prime mixes to within
a polylogarithmic factor on every input (Section~\ref{sec:energy}).
\item It runs the full pipeline on adversarial families and finds that low mixing
coincides with easy arithmetic structure, while hard families mix and run at the
sub-solver cost (Section~\ref{sec:exp}).
\end{itemize}
Section~\ref{sec:barrier} states what the experiments suggest the real
obstruction is. We do not claim a faster worst-case algorithm.

\subsection{A worked example of mixing}
Let the support be $A=\{3,5,8,13\}$ and let $w=2$, so the representations are the
$2$-subsets. Their sums are $8,11,16,13,18,21$. Modulo $p=7$ these are
$1,4,2,6,4,0$, so $s_p=5$ classes are hit and the coverage is $5/7\approx0.71$.
Modulo $p=5$ they are $3,1,1,3,3,1$, so only $s_p=2$ classes are hit and the
coverage is $0.4$. A random prime rarely looks like $p=5$; the point of this note
is to say how rarely, and what poor coverage implies.

\section{Preliminaries and the mixing condition}\label{sec:mix}
Let $A=\{a_1,\dots,a_m\}\subseteq\mathbb Z\setminus\{0\}$ and let $w$ with
$1\le w\le m-1$. In the balanced regime $m=n/2$ is the support of a solution and
$w=n/4$. Write $[m]=\{1,\dots,m\}$, and
\[
\mathcal Y=\{y\subseteq[m]:|y|=w\},\qquad
\Sigma(y)=\sum_{i\in y}a_i,\qquad Y=\binom mw ,
\]
so $\mathcal Y$ is the set of \emph{representations}. For a prime $p$ let
\[
c_r=\#\{y\in\mathcal Y:\Sigma(y)\equiv r\!\!\pmod p\},\qquad
E_p=\sum_r c_r^2 ,
\]
the \emph{collision energy}, and let $s_p=\#\{r:c_r>0\}$ be the number of residue
classes hit. The filter keeps the sums modulo $p$; a representation survives in a
random class with probability $s_p/p$, so the technique needs $s_p$ close to $p$.
We call $s_p/p$ the \emph{coverage}. Since $\sum_r c_r=Y$ and there are $s_p$
nonzero terms, Cauchy--Schwarz gives
\begin{equation}\label{eq:cs}
E_p\ge Y^2/s_p .
\end{equation}
The \emph{exact} collision energy is $E_\infty=\sum_v d_v^2$, where
$d_v=\#\{y:\Sigma(y)=v\}$.

\section{Mixing forces concentration}\label{sec:lemmas}
The first lemma says that poor mixing is a statement about the values, not merely
about their weight-$w$ sums. All proofs are in Appendix~\ref{app:proofs}.

\begin{lemma}[Concentration]\label{lem:conc}
$|(A\bmod p)-(A\bmod p)|\le s_p^2$, and hence $|\{a_i\bmod p\}|\le s_p^2+1$.
\end{lemma}

Poor mixing therefore forces additive structure. Structure is not the same as a
small sumset: $A=\{p\,2^0,\dots,p\,2^{m-1}\}$ has $s_p=1$ yet $|S(A)|=2^m$; it is
structured, and easy. The single-residue case is exact.

\begin{lemma}[Contraction]\label{lem:contract}
If $p\mid a_i-\rho$ for all $i$, then every representation sum satisfies
$\Sigma(y)\equiv w\rho\pmod p$. Writing $a_i=\rho+p\,b_i$, every subset sum is
$|y|\rho+p\,\Sigma_b(y)$, so $p$ can be stripped from the values and the instance
recursed.
\end{lemma}

For larger $s_p$, the analogue is that $A\bmod p$ lies in a short generalized
arithmetic progression (Kneser's theorem); we do not pursue that step here.

\section{Average collision energy over primes}\label{sec:energy}
The next result holds for \emph{every} input, on average over primes.

\begin{lemma}[Average energy]\label{lem:avg}
Let $B=\max_i|a_i|$, let $\kappa=\lfloor\log(2wB)/\log P\rfloor$, and let
$\pi=\#\{p\text{ prime}:P\le p\le2P\}$. Then
\[
\frac1\pi\sum_{p\in[P,2P]}E_p \;\le\; E_\infty+\frac{Y^2\kappa}{\pi}.
\]
\end{lemma}

Take $P\approx Y$, so $\pi\approx Y/\ln Y$, and $B\le2^{\beta n}$ with
$m,w=\Theta(n)$; then $\kappa\le2\beta+o(1)$. If few exact collisions occur,
i.e.\ $E_\infty\le Y\,\mathrm{polylog}(Y)$, Lemma~\ref{lem:avg} gives an average
of $\Oh(Y\log Y)$. Markov's inequality controls the tail: for all but an
$\Oh(1/\log Y)$ fraction of primes, $E_p=\Oh(Y\log^3Y)$, so by \eqref{eq:cs}
\begin{equation}\label{eq:cov}
s_p\;\ge\;\frac{Y^2}{E_p}\;=\;\frac{Y}{\mathrm{polylog}(Y)} .
\end{equation}
A representation is therefore found after $\mathrm{polylog}(Y)$ retries, and at
least half of the primes mix to coverage $Y/\mathrm{polylog}$.

Table~\ref{tab:energy} checks the bound directly. The average energy is about
$2Y$, the birthday value, well below the bound and below $Y\log Y$; and even the
\emph{pointwise minimum} coverage over all primes stays at $0.27$--$0.37$, not
only the average.

\begin{table}[t]
\centering
\caption{Average energy and coverage over every prime $p\in[Y,2Y]$, for a
planted balanced support with sparse values and no exact collisions.}
\label{tab:energy}
\begin{tabular}{rrrrrr}
\toprule
$n$ & $Y$ & primes & avg.\ $E_p$ & bound $Y^2\kappa/\pi$ & min.\ coverage\\
\midrule
$16$ & $70$  & $15$ & $113$ & $653$  & $0.371$\\
$20$ & $252$ & $43$ & $422$ & $2954$ & $0.325$\\
$22$ & $462$ & $68$ & $778$ & $6278$ & $0.270$\\
\bottomrule
\end{tabular}
\end{table}

\section{The full pipeline on adversarial families}\label{sec:exp}
Lemma~\ref{lem:avg} concerns all weight-$w$ sums; the algorithm needs a
representation of the actual solution. We therefore ran the complete pipeline:
enumerate weight-$n/8$ subsets of each half and match by residue (the balanced
sub-solver, $2\binom{n/2}{n/8}=2^{0.4057n}$ work per residue), filter by a random
prime, and look for a disjoint pair of weight-$n/4$ pieces with exact sum $b$.
Table~\ref{tab:adv} reports, for planted balanced solutions, the fraction solved
within $16$ residues, the coverage, and the enumeration work.

\begin{table}[t]
\centering
\caption{Full pipeline on adversarial families at $n=40$. Coverage is the
fraction of residues hit by the planted solution's representations; work is the
mean enumeration, baseline $2^{0.4057n}$.}
\label{tab:adv}
\begin{tabular}{lrrr}
\toprule
family & success & coverage & work\\
\midrule
random $\beta=1$        & $1.00$ & $0.95$ & $2^{0.4057n}$\\
random $\beta=1.5$      & $1.00$ & $0.95$ & $2^{0.4057n}$\\
two-scale $q b_i+r_i$   & $1.00$ & $0.94$ & $2^{0.4057n}$\\
rank-$2$ generalized AP & $1.00$ & $0.91$ & $2^{0.4057n}$\\
$q$-multiple            & $1.00$ & $0.95$ & $2^{0.4057n}$\\
geometric (superincr.)  & $0.17$ & $0.94$ & $>2^{0.4057n}$\\
arithmetic (large step) & $0.33$ & $0.003$ & $>2^{0.4057n}$\\
\bottomrule
\end{tabular}
\end{table}

The hard families---large random values, two-scale weights $qb_i+r_i$, rank-$2$
generalized progressions, near-multiples---all mix at $0.91$--$0.95$ and are
solved at the sub-solver cost, up to a small constant. The only low-coverage
family is the large-step arithmetic progression (coverage $0.003$), which is
additively structured and easy. Geometric (superincreasing) instances are
high-coverage but the balanced enumeration happens to fail on them, and they are
solved greedily in any case. No hard, poorly-mixing family was found.

\section{Where the barrier lies}\label{sec:barrier}
Combining the lemmas and the experiments: on every input a random prime mixes to
within a polylog (Lemma~\ref{lem:avg}), and in the tests the only low-coverage
families are structured and easy. If mixing were the obstruction, the pipeline
would already beat $2^{n/2}$ using the balanced sub-solver, since
$0.4057<1/2$. It does not, so for $\{0,1\}$ the obstruction must lie elsewhere.
The experiments isolate two candidates.

\emph{(a) The balanced sub-solver.} The filtered class is enumerated by matching
weight-$n/8$ subsets of two fixed halves, which finds only representations that
are weight-balanced across the split. An adversary can force the solution's
representations off balance, and the recursion that replaces the balanced
sub-solver---the four-way decomposition of~\cite{hgj,bcj}---is where the
published $0.291n$ bound is won; without it the exponent is stuck at $0.4057$.

\emph{(b) Compatibility and exactness.} Two pieces may combine to an invalid
coefficient vector (a ``pseudo-solution''), and recovering a compatible pair with
an exact, not merely modular, sum is a sparse orthogonal-vectors step whose cost
must be kept list-linear.

The open problem suggested by the experiments is therefore:
\begin{quote}
\itshape
For every $\{0,1\}$ instance, either a solution has a representation that the
balanced sub-solver enumerates and the compatibility certificate recovers, or the
instance is solved below $2^{n/2}$ by other means; and the filtered class is
enumerable and searchable in $2^{(1/2-\varepsilon)n}$ for a constant
$\varepsilon>0$.
\end{quote}

\emph{Remark on a recent preprint.} A 2025 preprint~\cite{salas} claims a
guaranteed $2^{n/2-\varepsilon}$ via a ``controlled aliasing'' that replaces one
element by another inside a half's partial sums and keeps one canonical subset
per aliased sum. The mechanism is unsound: in the half $\{1,2\}$, aliasing
$2\to1$ gives the canonical subset $\{1,2\}$ a stored sum of $2$ although its
true sum is $3$, so the merge can return a false positive and the subset
$\{2\}$ is lost. The claimed theorem is deferred to a companion paper and is not
established.

\section{Conclusion}
Mixing in the representation technique obeys an elementary, input-independent
bound: on average over primes the collision energy is $\Oh(Y\log Y)$, and
pointwise coverage stayed above $0.27$ on every input we tested. In the full
pipeline the hard, unstructured families mix at $0.9$ and run at the balanced
sub-solver cost, while low mixing coincides with arithmetic structure and easy
inputs. For the coefficient set $\{0,1\}$ this points to the sub-solver and the
compatibility step, not mixing, as the obstruction. A genuinely hard,
poorly-mixing family, if one exists, was not found; finding one, or proving none
exists, is the next step.

\appendices
\section{Proofs}\label{app:proofs}
\subsection{Proof of Lemma~\ref{lem:conc}}
Fix $i,j$. Since $1\le w\le m-1$ and $m\ge2$, the set $[m]\setminus\{i,j\}$ has at
least $m-2\ge w-1$ elements, so there is $S\subseteq[m]\setminus\{i,j\}$ with
$|S|=w-1$. The sets $S\cup\{i\}$ and $S\cup\{j\}$ both have size $w$, so
$u=a_i+\Sigma(S)$ and $v=a_j+\Sigma(S)$ are both values of weight-$w$ sums. Their
difference is $u-v=a_i-a_j$. Letting $i,j$ vary, every value difference modulo
$p$ is a difference of two hit residues; the hit residues number $s_p$, so the
difference set has at most $s_p^2$ elements. Fixing one index as a base gives
$|\{a_i\bmod p\}|\le s_p^2+1$. \hfill$\square$

\subsection{Proof of Lemma~\ref{lem:contract}}
If $a_i\equiv\rho\pmod p$ for all $i$, then
$\Sigma(y)=\sum_{i\in y}a_i\equiv|y|\rho=w\rho\pmod p$. Writing
$a_i=\rho+p\,b_i$ with $b_i=(a_i-\rho)/p\in\mathbb Z$, we get
$\Sigma(y)=|y|\rho+p\,\Sigma_b(y)$, so the target determines $p\,\Sigma_b(y)$ given
$|y|$. \hfill$\square$

\subsection{Proof of Lemma~\ref{lem:avg}}
For each prime $p$, $E_p=\sum_r c_r^2=\#\{(y,y'):\Sigma(y)\equiv\Sigma(y')\pmod p\}$,
so $\sum_pE_p=\sum_{y,y'}n_{y,y'}$ with
$n_{y,y'}=\#\{p\in[P,2P]:p\mid\Sigma(y)-\Sigma(y')\}$. If $\Sigma(y)=\Sigma(y')$
then $n_{y,y'}=\pi$; there are $E_\infty$ such ordered pairs. Otherwise
$N=\Sigma(y)-\Sigma(y')\ne0$ and $|N|\le2wB$ because each value is at most $B$ in
absolute value and a weight-$w$ sum has at most $w$ terms; the number of primes in
$[P,2P]$ dividing $N$ is at most $\lfloor\log|N|/\log P\rfloor\le\kappa$. The
remaining $Y^2$ ordered pairs each contribute at most $\kappa$. Summing and
dividing by $\pi$ gives the claim. \hfill$\square$

\begin{thebibliography}{9}
\bibitem{hs} E. Horowitz and S. Sahni, ``Computing partitions with applications
to the knapsack problem,'' \emph{Journal of the ACM}, vol.~21, no.~2,
pp.~277--292, 1974.
\bibitem{ss} R. Schroeppel and A. Shamir, ``A $T=O(2^{n/2})$, $S=O(2^{n/4})$
algorithm for certain NP-complete problems,'' \emph{SIAM Journal on Computing},
vol.~10, no.~3, pp.~456--464, 1981.
\bibitem{hgj} N. Howgrave-Graham and A. Joux, ``New generic algorithms for hard
knapsacks,'' in \emph{EUROCRYPT}, 2010; IACR ePrint 2010/189.
\bibitem{bcj} A. Becker, J.-S. Coron, and A. Joux, ``Improved generic algorithms
for hard knapsacks,'' in \emph{EUROCRYPT}, 2011; IACR ePrint 2011/474.
\bibitem{chen} X. Chen, Y. Jin, T. Randolph, and R. A. Servedio, ``Subset sum
in time $2^{n/2}/\mathrm{poly}(n)$,'' in \emph{RANDOM}, 2023; arXiv:2301.07134.
\bibitem{rw} T. Randolph and K. W\k{e}grzycki, ``Beating meet-in-the-middle for
subset balancing problems,'' arXiv:2511.10823, 2025.
\bibitem{akkn} P. Austrin, P. Kaski, M. Koivisto, and J. Nederlof, ``Dense subset
sum may be the hardest,'' in \emph{STACS}, 2016; arXiv:1508.06019.
\bibitem{salas} J. Salas, ``Beyond worst-case subset sum,'' arXiv:2503.20162,
2025.
\end{thebibliography}
\end{document}
